%! TEX root = master.tex
\chapter{4}{Real Integrals by Residues}

Some real integrals that resist every calculus technique become routine
once they are read as contour integrals. The simplest case is an integral of
a rational function of $\cos\theta$ and $\sin\theta$ over a full period, which
becomes an integral around the unit circle.

\section{Integrals over a period}\label{sec:trig-integrals}

%! TEX root = ../master.tex
\begin{theorem}[Trigonometric integrals]\label{thm:trig-integrals}
  Let $R(x, y)$ be a rational function with no poles on the circle
  $x^2 + y^2 = 1$. Then
  \[
    \int_0^{2\pi} R(\cos\theta, \sin\theta) \dd\theta
    = \oint_{\abs{z} = 1}
      R\left( \frac{z + z^{-1}}{2}, \frac{z - z^{-1}}{2i} \right)
      \frac{\dd z}{iz},
  \]
  and the right-hand side is $2\pi i$ times the sum of the residues of its
  integrand inside $\abs{z} = 1$.
  \thmsource{\cite[ch.~6]{saff-snider}}
\end{theorem}

\begin{proof}
  Parametrise the circle by $z = e^{i\theta}$, $0 \leq \theta \leq 2\pi$. Then
  $\dd z = iz \dd\theta$, $\cos\theta = \frac12(z + z^{-1})$ and
  $\sin\theta = \frac{1}{2i}(z - z^{-1})$, which turns the contour integral
  into the real one. The integrand on the right is rational in $z$, so it has
  finitely many poles, and \cref{thm:residue-theorem} applies.
\end{proof}


\begin{exercise}\label{ex:one-plus-quarter-cos}
  Evaluate $\displaystyle \int_0^{2\pi} \frac{\dd\theta}{1 + \frac14 \cos\theta}$.
\end{exercise}

\begin{solution}
  \emph{To the circle.} By \cref{thm:trig-integrals}, with $C$ the circle
  $\abs{z} = 1$,
  \[
    I = \oint_C \frac{1}{1 + \frac18\left(z + z^{-1}\right)} \frac{\dd z}{iz}
      = \oint_C \frac{8z}{z^2 + 8z + 1} \frac{\dd z}{iz}
      = \frac{8}{i} \oint_C \frac{\dd z}{z^2 + 8z + 1}.
  \]

  \emph{Poles.} The roots of $z^2 + 8z + 1$ are
  $z_{1,2} = -4 \pm \sqrt{15}$, both simple. Since $\sqrt{15} \approx 3.87$,
  $z_1 = -4 + \sqrt{15} \approx -0.13$ lies inside $C$ and
  $z_2 = -4 - \sqrt{15} \approx -7.87$ lies outside.

  \emph{Residue.} By \cref{thm:residue-at-pole},
  \[
    \Res\left( \frac{1}{(z - z_1)(z - z_2)}; z_1 \right)
    = \frac{1}{z_1 - z_2} = \frac{1}{2\sqrt{15}}.
  \]

  \emph{Total.}
  \[
    I = \frac{8}{i} \cdot 2\pi i \cdot \frac{1}{2\sqrt{15}}
      = \frac{8\pi}{\sqrt{15}}. \qedhere
  \]
\end{solution}

\begin{exercise}\label{ex:one-plus-a-cos}
  Show that for real $a$ with $0 < \abs{a} < 1$,
  \[
    \int_0^{2\pi} \frac{\dd\theta}{1 + a\cos\theta} = \frac{2\pi}{\sqrt{1 - a^2}}.
  \]
\end{exercise}

\begin{solution}
  As above,
  \[
    \int_0^{2\pi} \frac{\dd\theta}{1 + a\cos\theta}
    = \frac{2}{ia} \oint_{\abs{z} = 1} \frac{\dd z}{z^2 + \frac{2}{a} z + 1}.
  \]
  The roots $z_{1,2} = \bigl(-1 \pm \sqrt{1 - a^2}\bigr)/a$ are real with
  product $1$ and $\abs{z_1} < \abs{z_2}$, so exactly one of them, $z_1$, lies
  inside the circle. Its residue is
  $1/(z_1 - z_2) = a / \bigl(2\sqrt{1 - a^2}\bigr)$, and so
  \[
    \int_0^{2\pi} \frac{\dd\theta}{1 + a\cos\theta}
    = \frac{2}{ia} \cdot 2\pi i \cdot \frac{a}{2\sqrt{1 - a^2}}
    = \frac{2\pi}{\sqrt{1 - a^2}}.
  \]
  With $a = \frac14$ this gives $2\pi / \sqrt{15/16} = 8\pi / \sqrt{15}$,
  agreeing with \cref{ex:one-plus-quarter-cos}.
\end{solution}

\begin{remark}
  As $a \to \pm 1$ the integral blows up: one pole moves onto the unit circle,
  where $1 + a\cos\theta$ vanishes.
\end{remark}
